\chapter{Newton Polygon}

\section{Definitions}

In this section we will determine a way to check for zeroes of a power series via the construction
of Newton polygons. We consider power series over a non-archimedean valued field $R$ where we denote
the \emph{additive} valuation by $v$ and the associated multiplicative valuation by $\lvert \cdot \rvert$.
The typical example would be $\Q_p$.

These are formally given by the following definition.

\begin{definition}
    The \emph{Newton polygon} of $f = \sum a_i x^i$ is the lower boundary of the convex hull
    of the set of points $(i,\nu (a_i))$, where we ignore points with infinite $p$-adic valuation.
    Here by convex hull we mean the smallest convex polygon that contains all points.
\end{definition}

Given a Newton polygon we can give the following information.

\begin{definition}
    \begin{enumerate}
        \item The \emph{slopes} of the line segments appearing in the Newton polygon are called the
        slopes of $f(x)$.
        \item The \emph{length} of a slope is the length of the projection of the line segment with
        that slope onto the $x$-axis.
        \item The \emph{breaks} are the values of $i$ which are vertices of the Newton polygon.
    \end{enumerate}
\end{definition}

It turns out when working with Newton polygons the slopes and lengths are the information we care about;
and moreover given a starting point using this we can reconstruct all breaks. This provides motivation
for the definition we give in Lean.

\begin{definition}
  A Newton polygon is the following information:
  \begin{itemize}
    \item a natural number $n$, determining the number of slopes;
    \item a (potentially infinite) list of slopes that are strictly increasing except for the final
      slope (which may only be equal to the previous);
    \item a (potentially infinite) list of lengths; and
    \item a starting point.
  \end{itemize}
\end{definition}

It is clear that given a Newton polygon, in the formal sense, we can construct the following information;
however the reverse also holds by using the comment above the definition and the fact the slopes are
(strictly) increasing give the convexity condition.
%% Perhaps I could elaborate on this further?

We now want to construct a way to obtain this information given a power series, $f = \sum_{i=0}^\infty a_n x^i \in R[x]$.
Towards this we consider the following algorithm:

\begin{enumerate}
    \item Plot the points $\{(i,\nu (a_i) )\}$.
    \item Start with the vertical half-line down from the point with the smallest $i$.
    \item Rotate the line until one of the following happens:
    \subitem The line simultaneously hits infinitely many of the points we have plotted. In this
        case don't break the line and the polygon is complete.
    \subitem The line reaches a position where it can be rotated no further without leaving behind
        some points. In this case, the half lines forms the final segment of the polygon.
    \subitem The line hits a finite number of points. In this case, break at the last point that was
        hit and repeat the procedure.
\end{enumerate}

It should be clear that by construction this constructs the lower convex hull of the set of points
$\{(i,\nu (a_i) )\}$ and taking the slopes and lengths of the obtained segments we get our corresponding
definition in Lean.

As mentioned, the study of Newton polygons is deeply connected to studying the zeroes of $f(x)$,
thus we will assume $f(0) \neq 0$. That is, we factor out any powers of $x$ dividing $f$, which
means our initial vertical line starts on the $y$-axis.

Moreover, we will also assume $a_0 = 1$. This is because if we factor
$f(x) = \sum_{i = 0}^\infty a_i x^i$ as $f(x) = a_0 \cdot g(x)$, where
$g(x) = 1 + \sum_{i=1}^\infty \frac{a_i}{a_0}x^i$, it is not hard to see that the resulting Newton
polygon of $g(x)$ is the Newton polygon of $f(x)$ translated down by $\nu (a_0)$.
This fact follows from observing $\nu (\frac{a_i}{a_0}) = \nu(a_i) - \nu (a_0)$.

\section{Properties}

In this section we will consider power series over $\Q_p$ - however, we can easily generalise to
complete non-archimedean valued fields.

The first thing we can use the Newton polygon for is to obtain the radius of convergence of the
power series.

\begin{theorem} \label{RadofCon}
    Let $m$ be the supremum of the slopes appearing in the Newton polygon of a series
    $f(x) = 1 + a_1 x + a_2 x^2 + \cdots$. Then the radius of convergence of $f$ is $p^m.$
\end{theorem}



\begin{proposition}\label{NP}
    Let $f(x) = 1 + a_1 x + a_2 x^2 + \cdots$ be a power series. Let $m_1, m_2, \dots, m_k$ be the
    first $k$ slopes of the Newton polygon of $f(x)$, and assume that $f(x)$ converges on the
    closed ball of radius $c = p^{m_k}$ (from previous section, this means it is a restricted power
    series for parameter $c$ if $K$ is complete). Let $N$ be the $x$-coordinate of the right
    endpoint of the $k$-th segment of the Newton polygon. Then, there exists a polynomial
    $g(x)$ of degree $N$ and a power series $h(x)$ such that
    \begin{enumerate}
        \item $f(x) = g(x) h(x)$,
        \item $\| f(x) - g(x) \|_c < 1$,
        \item $h(x)$ converges on the closed ball of radius $c$,
        \item $\| h(x) - 1\|_c <1$, and
        \item the Newton polygon of $g(x)$ is the same as the portion of the Newton polygon of
        $f(x)$ contained in the region $0 \leq x \leq N$.
    \end{enumerate}
\end{proposition}

This then allows us to give our connection to zeroes of the power series.

\begin{corollary}
    Let $f(x) = 1 + a_1 x + a_2 x^2 + \cdots$ be a power series which converges on the closed ball
    of radius $c = p^m$. Let $m_1,\dots, m_k$ be the slopes of the Newton polygon of $f(x)$ which
    are less than or equal to $m$, and let $i_1,\dots,i_k$ be their lengths. Then, for each $j$, the
    power series $f(x)$ has $i_j$ zeros with absolute value $p^{m_j}$, and there are no other zeros
    in the closed ball of radius $p^m.$
\end{corollary}

\begin{proof}[Proof of proposition~\ref{NP}]
    We prove this by induction on $k$:\\
    If $k=1$, write the break point as $(N, m_1 N)$. We know all points are on or above the line
    $y = m x$, and the ones where $j > N$ are strictly above it. This translates to
    \begin{itemize}
        \item $| a_j | (p^j)^{m_1} \leq 1$ for all $j$,
        \item $| a_N | (p^N)^{m_1} = 1$, and
        \item $| a_j | (p^{m_1})^j$ < 1 if $j > N.$
    \end{itemize}
    This says that $\|f(x) \|_{p^{m_1}} = 1$ and that the maximum is last realised at the degree
    $N$. Thus, we can apply the Weierstrass preparation theorem to conclude that there is a
    polynomial of degree $N$ and a power series $h(x)$ satisfying
    \[
    \| h(x) - 1 \|_{p^m_1} < 1 \quad \|f(x) - g(x)\|_{p^m_1} < \|f(x) \|_{p^m_1} \quad \text{and} \quad f(x) = g(x) h(x).
    \]
    Combining the middle equation with $\| f(x) \|_{p^{m_1}} < 1$ means
    $\| f(x) - g(x) \|_{p^{m_1}} < 1$.
    Now, using facts about Newton polygons of polynomials, it follows that the zeros of $f(x)$ in
    the closed ball of radius $p^m$ coincide with the those of $g(x)$. That is, we must have the
    Newton polygon of $g(x)$ coincides with the Newton polygon of $f(x)$ on $0 \leq x \leq N$.
    Finally, we have $f(x)$ converges on the closed ball of radius $p^{m_1}$, by Theorem
    \ref{RadofCon}, and thus so does $h(x).$
    \\
    By induction, assume we have $f(x) = g_{k-1}(x)h_{k-1}(x)$ for a polynomial $g_{k-1}(x)$ and
    power series $h_{k-1}(x)$, such that the Newton polygon of $g_{k-1}(x)$ coincides with the
    first $k-1$ segments of the polygon for $f(x)$ and $h_{k-1}(x)$ has no zeroes in the closed ball
    of radius $p^{m_{k-1}}$.
    \\
    Let us now consider the $k$-th segment. By definition, the $k$-th segment ending at $x = N$
    means that for any $i > N$, $(i, \nu_p (a_i))$ lies above the line of slope $m_k$ through the
    point $(N, \nu_p (a_N))$. This means that, writing $c = p^{m_k}$, $\| f(x) \|_{p^{c}} =
    \lvert a_N \rvert c^N$ and $\lvert a_n \rvert c^N > \lvert a_i \rvert c^i $ for any $i > N$.
    Thus, we can apply the $p$-adic Weierstrass preparation theorem to get a polynomial $g(x)$.
    It is then easy to see $g_1 (x) \mid g(x)$, and that its Newton polygon coincides with the
    relevant portion of the Newton polygon of $f(x).$
\end{proof}


%% I need to re-go over this section and the proofs... plus I need to check the difference
%% between BGR Weierstrass prep and Gouvea's -- i.e. when I apply it do I need extra justification...
